\uuid{Ac6y}
\exo7id{3781}
\titre{exo7 3781}
\auteur{quercia}
\organisation{exo7}
\datecreate{2010-03-11}
\isIndication{false}
\isCorrection{false}
\chapitre{Espace euclidien, espace normé}
\sousChapitre{Endomorphismes auto-adjoints}
\cpge{mpsi-mp2i.calcul-matriciel-systemes.e, pcsi.determinants.c, ptsi.determinants.a, mp-mpi.endomorphismes-euclidien.b, pc.endomorphismes-euclidiens.b, psi.endomorphismes-euclidien.b, pt.isometries-euclidien.b, pc.complements-espaces-vectoriels.c, psi.complements-espaces-vectoriels.c, pt.complements-algebre-lineaire.c, bcpst2.produit-scalaire.c, bcpst2.produit-scalaire.a}
\module{Algèbre}
\niveau{L2}
\difficulte{}

\contenu{
\question{
Soit $A = (a_{ij}) \in \mathcal{M}_n(\R)$ symétrique de valeurs propres
$\lambda_1,\dots,\lambda_n$.
Montrer que $\sum_{i,j} a_{ij}^2 = \sum_i \lambda_i^2$.
}
}
